\section{Limitations and Conclusion}
\label{sec:conclusion}
\textbf{Limitations.} SPECTRA-Bench stops at schematic level: no layout parasitics, and passives use a first-order
poly-resistor model and ideal capacitors. The certified tier derives from our own templates; the independently written designer tier mitigates
but does not remove this bias. Compute (4 CPU cores) limited each task to one to three seeds, so per-task
estimates are coarse and the confidence intervals are wide; for the same reason we dropped a Gen+CMA-ES arm, an
all-corners PVT strategy, a template oracle, a cross-backbone arm, and infeasible-specification controls, which we built
incorrectly and withdrew. All reported LLM results use GPT-6.1-sol, so this study does not establish cross-backbone
generality. Reviewer-side checks of this work were performed by models of the same family and are therefore provisional.

\textbf{Conclusion.} Treating the LLM as a designer of \emph{intent} rather than of numbers---writing executable
$g_m/I_D$ sizing programs that predict their own behaviour---and closing the loop with device-level diagnosis and
intent-space, worst-corner-aware refinement gives a reproducible end-to-end protocol for tight, many-specification
and PVT-robust tasks. On the frozen external SIZ7g matrix SPECTRA reaches TT/PVT success
$\SizSpectraTT{}/\SizSpectraPVT{}$, while AnaFlow reaches $\SizAnaFlowTT{}/\SizAnaFlowPVT{}$.
The matched Xyce sensitivity shows lower PVT violation on both public topologies but retains a TT deficit on HoiLee,
and the SGAnalog probe exposes an explicit Xyce model-compatibility boundary. We release SPECTRA, SPECTRA-Bench,
all run logs and the gate checker so future work can extend the measured strata without conflating capability audits
with ranking results. The current receipts support a fixed-topology benchmark claim; they do not support a universal
all-circuit, all-PDK, all-simulator SOTA claim.
